\chapter{Regularization in Deep Learning}

\section{Dropout and Noise Injection}

Deep neural networks are highly expressive and can easily overfit training data. Regularization techniques are therefore essential for improving generalization.

\medskip

\noindent
\textbf{Dropout.}  
Dropout is a stochastic regularization method in which units are randomly deactivated during training.

\medskip

\noindent
\textbf{Mechanism.}  
For each unit, a binary mask is sampled:
\[
m_i \sim \text{Bernoulli}(p),
\]
and activations are modified as
\[
\tilde{z}_i = m_i z_i.
\]

\medskip

\noindent
\textbf{Interpretation.}
\begin{itemize}
    \item Each iteration trains a different subnetwork
    \item Prevents co-adaptation of features
\end{itemize}

\medskip

\noindent
\textbf{Inference.}  
At test time, activations are scaled to account for dropout:
\[
z_i \leftarrow p \, z_i.
\]

\medskip

\noindent
\textbf{Effect.}
\begin{itemize}
    \item Improves robustness
    \item Reduces overfitting
\end{itemize}

\medskip

\noindent
\textbf{Noise injection.}  
More generally, noise can be added to inputs, activations, or parameters:
\[
z \leftarrow z + \varepsilon,
\]
where $\varepsilon$ is random noise.

\medskip

\noindent
\textbf{Insight.}  
Noise encourages the model to learn stable features that are insensitive to small perturbations.

\section{Weight Decay}

Weight decay is one of the most widely used regularization techniques in deep learning.

\medskip

\noindent
\textbf{Objective.}
\[
\min_\theta \; \hat{R}_n(\theta) + \lambda \|\theta\|_2^2.
\]

\medskip

\noindent
\textbf{Interpretation.}
\begin{itemize}
    \item Penalizes large parameter values
    \item Encourages simpler models
\end{itemize}

\medskip

\noindent
\textbf{Effect on optimization.}  
The update rule becomes
\[
\theta_{t+1} = (1 - \eta \lambda)\theta_t - \eta \nabla \hat{R}_n(\theta_t).
\]

\medskip

\noindent
\textbf{Insight.}  
Parameters are gradually shrunk toward zero during training.

\medskip

\noindent
\textbf{Connection to stability.}  
Smaller weights often lead to smoother functions and improved generalization.

\medskip

\noindent
\textbf{Bayesian interpretation.}  
Weight decay corresponds to a Gaussian prior on parameters.

\section{Data Augmentation}

Data augmentation improves generalization by artificially increasing the diversity of the training data.

\medskip

\noindent
\textbf{Basic idea.}  
Generate new training examples by applying transformations:
\[
(x, y) \;\rightarrow\; (T(x), y),
\]
where $T$ preserves the label.

\medskip

\noindent
\textbf{Examples (vision).}
\begin{itemize}
    \item Rotation, translation, scaling
    \item Flipping, cropping
    \item Color perturbations
\end{itemize}

\medskip

\noindent
\textbf{Examples (other domains).}
\begin{itemize}
    \item Noise injection in inputs
    \item Synonym replacement in text
\end{itemize}

\medskip

\noindent
\textbf{Interpretation.}
\begin{itemize}
    \item Encodes invariances (e.g., rotation invariance)
    \item Expands effective dataset size
\end{itemize}

\medskip

\noindent
\textbf{Geometric view.}  
Augmentation enlarges the region of input space associated with each label.

\medskip

\noindent
\textbf{Insight.}  
Data augmentation acts as a form of regularization by constraining the function to behave consistently under transformations.

\section{Implicit Bias of Optimization}

Regularization in deep learning is not limited to explicit techniques. Optimization itself introduces a form of implicit bias.

\medskip

\noindent
\textbf{Overparameterization.}  
Deep networks often have more parameters than data points, leading to many solutions with zero training error.

\medskip

\noindent
\textbf{Key question.}  
Which solution does the optimization algorithm select?

\medskip

\noindent
\textbf{Implicit bias.}  
Gradient-based methods tend to favor solutions with particular properties, such as:
\begin{itemize}
    \item small norms,
    \item smooth functions,
    \item certain structural patterns.
\end{itemize}

\medskip

\noindent
\textbf{Example (linear models).}  
Gradient descent converges to the minimum-norm solution among all interpolating solutions.

\medskip

\noindent
\textbf{Deep learning perspective.}
\begin{itemize}
    \item SGD noise influences the trajectory of optimization
    \item Different optimizers lead to different solutions
\end{itemize}

\medskip

\noindent
\textbf{Insight.}  
The training process itself acts as a regularizer.

\medskip

\noindent
\textbf{Interaction with explicit regularization.}  
Implicit and explicit regularization often work together to shape the learned model.

\section{A Unifying Perspective}

Regularization in deep learning takes multiple forms:

\begin{itemize}
    \item \textbf{Explicit:} weight decay, dropout
    \item \textbf{Data-driven:} augmentation
    \item \textbf{Implicit:} optimization dynamics
\end{itemize}

\medskip

\noindent
\textbf{Core objective.}
\[
\text{Fit data} \quad + \quad \text{promote generalization}.
\]

\medskip

\noindent
\textbf{Key insight.}  
Generalization is shaped not only by the model, but by how it is trained and what data it sees.

\section{Outlook}

This chapter examined regularization techniques specific to deep learning:
\begin{itemize}
    \item dropout and noise injection,
    \item weight decay,
    \item data augmentation,
    \item implicit bias of optimization.
\end{itemize}

\medskip

In the next chapter, we study residual networks and skip connections, which play a key role in enabling very deep architectures.

\medskip

\noindent
\textbf{Guiding principle.}  
Effective regularization arises from combining model structure, data transformations, and optimization dynamics.